\chapter{Implementación y validación}

\section{Tecnologías principales}

La fase 1 está implementada en Python y usa \texttt{uv} para ejecución y pruebas.
El paquete principal es \texttt{decisionlab}. La infraestructura compartida se
encapsula en servicios reutilizables para evitar dependencias directas con la
fase 2.

\begin{table}[H]
\begin{center}
\begin{tabular}{|l|p{9cm}|} \hline
\textbf{Tecnología} & \textbf{Uso en la fase 1} \\ \hline
Python & Implementación del pipeline y modelos generados. \\ \hline
PostgreSQL & Metadatos, ejecuciones, modelos y ciclo de vida de memoria. \\ \hline
MinIO & Almacenamiento de artefactos y estado de pipeline. \\ \hline
Neo4j & Grafo de conocimiento científico. \\ \hline
Qdrant & Recuperación densa y BM25 de hechos de memoria. \\ \hline
Embeddings & Similitud semántica y enlazado aproximado de entidades. \\ \hline
\end{tabular}
\caption{Tecnologías usadas por la fase 1.}
\end{center}
\end{table}

\section{Bucle de herramientas}

Los agentes se ejecutan mediante un bucle genérico de herramientas. El modelo
propone llamadas, el dispatcher las ejecuta, y los resultados vuelven al modelo.
Las llamadas se registran para auditoría.

\begin{verbatim}
mensaje + herramientas
       |
       v
modelo propone tool_use
       |
       v
dispatcher ejecuta herramientas
       |
       v
resultados + traza
       |
       v
modelo continua o termina
\end{verbatim}

Las herramientas se diseñan con entradas y salidas estrechas. Esto reduce
ambigüedad y permite que el Router mantenga control sobre persistencia,
revisión y memoria.

\section{Paralelismo controlado}

Formalizer, Reasoner y Builder pueden procesar varios paradigmas o
formulaciones en paralelo. El sistema recoge fallos por subagente para que un
error no oculte salidas válidas de otros elementos.

\begin{verbatim}
Paradigma A -> Subagente A -> resultado A
Paradigma B -> Subagente B -> resultado B
Paradigma C -> Subagente C -> error controlado
\end{verbatim}

Este diseño reduce tiempo de ejecución sin convertir el pipeline en una caja
negra. Los artefactos siguen teniendo rutas estables y trazas por ejecución.

\section{Contrato del modelo generado}

El modelo generado debe ser autocontenido. No hereda de una clase base del
simulador; implementa métodos con nombres y comportamiento esperados.

\begin{lstlisting}[language=Python]
class SomeDecisionModel:
    def decide(self, perception: dict) -> Action:
        ...

    def update(self, action, reward, new_perception) -> None:
        ...

    def get_state(self) -> dict:
        ...
\end{lstlisting}

La fase 2 ejecuta el modelo en el siguiente orden:

\begin{lstlisting}[language=Python]
pre_state = model.get_state()
perception = env.build_perception(agent)
action = model.decide(perception)
reward, result = env.apply(action)
new_perception = env.build_perception(agent)
model.update(action, reward, new_perception)
\end{lstlisting}

Separar \texttt{decide} y \texttt{update} evita que el modelo aprenda antes de
conocer recompensa y nueva percepción. También facilita reproducir y depurar una
decisión concreta.

\section{Niveles de validación}

La validación del sistema no depende de un único mecanismo.

\begin{itemize}
\item Validación estructural: schemas, Pydantic y JSON esperado.
\item Validación semántica: prompts con criterios explícitos para Reasoner y
Builder.
\item Revisión humana: aceptación o reejecución de artefactos relevantes.
\item Validación ejecutable: tests generados y ejecutados con \texttt{pytest}.
\item Persistencia y trazabilidad: artefactos, estado y \texttt{trace.jsonl}.
\end{itemize}

\begin{figure}[H]
\centering
\begin{verbatim}
Texto cientifico
  -> revision humana
  -> especificacion estructurada
  -> validacion semantica
  -> codigo
  -> pytest
  -> registro del modelo
\end{verbatim}
\caption{Validación acumulativa del resultado.}
\end{figure}

\section{Degradación controlada}

El sistema distingue infraestructura obligatoria y opcional. PostgreSQL y MinIO
son necesarios para persistir ejecuciones y artefactos. Neo4j, Qdrant,
embeddings y reranking pueden faltar; en ese caso, el pipeline sigue generando
modelos, pero sin memoria avanzada.

Esta decisión mejora la robustez local y facilita desarrollo. La contrapartida
es que la calidad de las decisiones de los agentes puede disminuir cuando no
existe recuperación de conocimiento.

\section{Trazabilidad}

Cada ejecución mantiene estado, artefactos y trazas. Esto es esencial en un
sistema con modelos de lenguaje: no basta con conocer la salida final, también
hay que poder reconstruir qué fuentes, herramientas y decisiones llevaron a
ella.

La trazabilidad permite:

\begin{itemize}
\item auditar informes científicos;
\item comparar formulaciones aprobadas y rechazadas;
\item depurar especificaciones no implementables;
\item reproducir fallos de Builder;
\item alimentar la memoria solo con conocimiento aceptado.
\end{itemize}
